% !TEX TS-program = lualatex+se+debug
% !TEX encoding = UTF-8 Unicode

\documentclass[14pt, twoside, a4paper]{extreport}
% -----------------------------------------------------------------------------------
% ------------------------- USTAWIENIA STRONY--------------------------
% -----------------------------------------------------------------------------------
\usepackage{geometry}
\geometry{
	left=10mm,
	right=10mm,
	top=15mm,
	bottom=15mm,
	marginparwidth=0mm,
	marginparsep=0mm
	}
\newcommand{\wstawplik}[1]{\input{../include/#1}}%skraca kod
\wstawplik{headerofficium}
% elementy zmienne co do headerofficium:
\graphicspath{{../img/}} % katalog obrazów
%\svgpath{{../include/}}
\grechangestaffsize{20}%rozmiar wyniku GregorioTEX
% koniec 
\gresetlastline{ragged}% justowanie do lewej
% debugowanie:
% %\iffalse %*** %\fi %***
% \tracingall przed problematycznym kodem - bardzo dużo informacji
% lualatex+se+debug - dodatkowe opcje: -file-line-error i -halt-on-error 
% \errorcontextlines
%\def\mydebug{1}
\begin{document}
\makeatletter
\grechangedim{interwordspacetext}{\unexpanded{\the\fontdimen2\font plus \the\gluestretch\gre@space@skip@interwordspacenotes minus \the\glueshrink\gre@space@skip@interwordspacenotes}}{fixed}%
\makeatother
\pagestyle{empty}
%\pagestyle{plain}	
\large
\begin{spacing}{0.7}
\begin{center}
	\section*{In Parasceve - lectio 5}
	%\small\textit{Fontevraud Gradual\\}
	\small\textit{Laudes festivae (1940) p.81}
	\vspace{5mm}
\end{center}	
\end{spacing}
%\greannotation{} 
\gresetinitiallines{1}\greannotation{Lectio 5}\grecommentary[2mm]{}\gabcsnippet{(f3)No(e)stis(f) qui(h) con(h)vén(h)tus(h) e(h)rat(h) ma(h)li(h)gnán(h)ti(h)um(h) Iu(h)dæ(h)ó(i)rum,(f) (;) et(h) quæ(h) mul(h)ti(h)tú(h)do(h) e(h)rat(h) o(h)pe(h)rán(h)ti(h)um(h) i(h)ni(e)qui(f)tá(g)tem.(f) (:) Quam(e) i(f)ni(h)qui(f)tá(h)tem?(i) (:) Qui(h)a(h) vo(h)lu(h)é(h)runt(h) oc(h)cí(h)de(h)re(h) Dó(h)mi(h)num(h) Ie(e)sum(f) Chri(g)stum.(f) (:) Tan(c)ta(ef) ó(f)pe(f)ra(f) bo(f)na,(f) in(f)quit,(f) o(f)stén(f)di(e) vo(f)bis:(f) (;) prop(e)ter(f) quod(h) ho(h)rum(h) me(h) vul(h)tis(h) oc(f)cí(h)de(h)re?(h) (:) Pér(h)tu(h)lit(h) om(h)nes(h) in(h)fír(h)mos(h) e(g)ó(f)rum,(f) (,) cu(f)rá(f)vit(f) om(f)nes(f) lán(f)gui(f)dos(f) e(e)ó(f)rum,(f) (;) præ(f)di(f)cá(f)vit(f) re(f)gnum(g) cæ(h)ló(i)rum,(h) (;) non(h) tá(h)cu(h)it(h) ví(h)ti(h)a(g) e(f)ó(gh)rum,(g) (;) ut(e) i(gh)psa(f) pó(f)ti(f)us(f) e(f)is(f) dis(g)pli(h)cé(i)rent,(h) (,) non(h) mé(h)di(h)cus,(h) a(h) quo(h) sa(e)na(f)bán(g)tur.(f) (:) His(c) óm(ef)ni(f)bus(f) cu(f)ra(f)ti(f)ó(f)ni(f)bus(f) e(f)ius(g) in(h)grá(g)ti,(f) (;) tam(e)quam(f) mul(h)ta(h) fe(h)bre(h) phre(h)né(i)ti(f)ci,(f) (;) in(h)sa(h)ni(h)én(h)tes(h) in(h) mé(h)di(h)cum,(h) (,) qui(h) vé(h)ne(h)rat(h) cu(h)rá(g)re(f) e(g)os,(gh) (;) ex(h)co(h)gi(h)ta(h)vé(h)runt(h) con(h)sí(h)li(h)um(h) per(h)dén(e)di(f) e(g)um:(f) (:) tam(e)quam(f) i(h)bi(h) vo(h)lén(h)tes(h) pro(h)bá(i)re,(f) (;) u(h)trum(h) ve(h)re(h) ho(h)mo(h) sit,(h) qui(h) mo(g)ri(f) pos(g)sit,(gh) (;) an(h) á(h)li(h)quid(h) su(h)per(h) hó(h)mi(h)nes(h) sit,(h) (,) et(h) mo(h)ri(h) se(h) non(e) per(f)mít(g)tat.(f) (:) Ver(e)bum(f) i(h)psó(h)rum(h) a(h)gnó(h)sci(h)mus(h) in(h) Sa(h)pi(h)én(h)ti(h)a(h) Sa(g)lo(f)mó(gh)nis:(g) (;) Mor(e)te(gh) tur(f)pís(f)si(f)ma,(e) ín(f)qui(f)unt,(f) (,) con(f)dem(f)né(g)mus(h) e(i)um.(h) (;) In(h)ter(h)ro(h)gé(g)mus(f) e(g)um:(gh) (;) e(h)rit(h) e(h)nim(h) re(h)spé(h)ctus(h) in(h) ser(h)mó(h)ni(h)bus(e) il(f)lí(g)us.(f) (:) Si(c) e(ef)nim(e) ve(fh)re(h) Fí(h)li(h)us(h) De(h)i(h) est,(i) (;) lí(h)be(e)ret(f) e(g)um.(f) (::)}\paragraphspace%
\tlumacz{Wiecie, że Zgromadzenie uczynili Żydowie przewrotni i wiecie jak wielką była liczba czyniących nieprawość? Tę iż chcieli zabić Pana Jezusa Chrystusa. Tyle dobrych dzieł uczyniłem dla was rzecze im: dla któregoż to chcecie mnie zabić? Przyjmował wszystkich ułomnych ich, uzdrawiał wszystkich chorych, opowiadał królestwo niebieskie, ale nie zamilczał występków ich. Nie podobało im się to lekarstwo, które im podawał i nie chcieli go mieć za lekarza, któryby ich uzdrowił. Nie pamiętni wszystkich dobrodziejstw Jego, podobni do szaleńców w gorączce, których drażni obecność lekarza, przychodzącego uzdrowić ich; uknuli spisek, aby Go zgubić. Zrobili to, jakby chcąc przekonać się, kim On był naprawdę: człowiekiem podległym śmierci, czy człowiekiem wyższym ponad resztę ludzi, nad którym śmierć władzy nie ma. W księdze Mądrości czytamy przepowiednię słów, które wtedy wyrzekli: Skażmy go, aby umarł śmiercią najhaniebniejszą, pytajmy go, abyśmy mogli mieć wzgląd na słowa Jego. Jeśli prawdziwie jest Synem Bożym, Bóg wybawi Go.}
\end{document}